\chapter{The Boot Path: Bootloader and PSRAM}
\label{ch:bootloader}

When the chip comes out of reset, a great deal has to happen before the first
line of Ada in your \texttt{Main} runs: the application has to be read out of
flash and placed in RAM, the instruction and data caches have to be pointed at
the right flash windows, the CPU has to be sped up from its sleepy reset clock,
the register windows have to be made sane for Ada's calling convention, and --
if the board has it -- the external PSRAM has to be brought up and mapped. This
chapter follows that path. It is the part of the system that is \emph{not}
Espressif's: the project replaced the last vendored binary, the 2nd-stage
bootloader, with its own -- written, like almost everything else here, in Ada.

\section{Three things run before \texttt{Main}}

Boot happens in three hand-offs (the boot figure in
Chapter~\ref{ch:anatomy}, Fig.~\ref{fig:boot}, sketches the same sequence):

\begin{enumerate}
\item \textbf{The mask ROM} (silicon, unchangeable) wakes on reset, reads the
  image at flash offset \texttt{16\#0\#}, and loads it into SRAM. That image is our
  \emph{2nd-stage bootloader}.
\item \textbf{Our 2nd-stage bootloader} reads the \emph{application} image (at
  flash offset \texttt{16\#1\_0000\#}), copies its RAM segments into place, maps its
  flash-resident code and data through the cache MMU, brings up PSRAM, and jumps
  to the application's entry point.
\item \textbf{The application's \texttt{start.S}} normalises the CPU state,
  configures the instruction cache, switches to the full clock, and calls into
  the Ada world (eventually the environment task that runs \texttt{Main}).
\end{enumerate}

\section{The 2nd-stage bootloader}\index{bootloader}

The loader is ZFP-style Ada -- no runtime under it -- in
\texttt{examples/common/bare/bootloader/} (the core is \texttt{boot\_main.adb}).
It is deliberately small. Its job, in order:

\begin{itemize}
\item \textbf{Silence the watchdogs.} The boot ROM arms the timer-group and RTC
  watchdogs; the loader unlocks them (the magic key \texttt{16\#50D8\_3AA1\#}) and
  clears them so they cannot reset the chip mid-load.
\item \textbf{Read and validate the app header.} At flash offset
  \texttt{16\#1\_0000\#} sits the image header: a \texttt{16\#E9\#} magic byte, a segment
  count, and the entry-point address. (This is exactly the format the
  \texttt{esp\_elf2image} host tool emits -- \S\ref{ch:build}.)
\item \textbf{Place each segment by where it wants to live.} For every segment
  the loader looks at its load address:
\begin{center}
\begin{tabular}{l l l}
\toprule
\textbf{Window} & \textbf{Range} & \textbf{Action} \\ \midrule
IRAM   & \texttt{16\#4037\_0000\#..16\#403E\_0000\#} & copy from flash into SRAM \\
DRAM   & \texttt{16\#3FC8\_0000\#..16\#3FD0\_0000\#} & copy from flash into SRAM \\
IROM   & \texttt{16\#4200\_0000\#..16\#4400\_0000\#} & map (execute-in-place via cache) \\
DROM   & \texttt{16\#3C00\_0000\#..16\#3E00\_0000\#} & map (read-in-place via cache) \\
\bottomrule
\end{tabular}
\end{center}
  RAM segments are copied byte-for-byte with the ROM SPI-flash read routine; the
  large code and read-only-data segments are left \emph{in} flash and reached
  through the cache.
\item \textbf{Map the flash windows through the cache MMU.} The loader programs
  the instruction-bus and data-bus MMUs (\texttt{Cache\_Ibus\_MMU\_Set} /
  \texttt{Cache\_Dbus\_MMU\_Set}, 64\,KB pages) so the IROM/DROM virtual
  addresses resolve to the right flash pages, then clears the cache ``shut'' bits
  to switch the buses on.
\item \textbf{Bring up PSRAM} (next section) and \textbf{jump} to the entry point.
\end{itemize}

The loader's own code and stack are placed high in SRAM, above the region the
application's segments occupy, so copying the app never overwrites the running
loader.

\section{PSRAM: octal bring-up}\index{PSRAM}

External PSRAM is the trickiest thing the bootloader does: bringing up octal PSRAM
reconfigures the very MSPI timing that flash execute-in-place depends on, so it has
to run from RAM, carefully, at exactly the right moment -- and the 2nd-stage
bootloader is that moment. The bring-up (a small C shim, \texttt{psram\_boot.c}) is
now almost entirely \emph{from source}; only the lowest-level OPI and cache pokes
still call masked-ROM helpers (\texttt{esp\_rom\_opiflash\_*}, \texttt{Cache\_*}).
Chapter~\ref{ch:psram} tells the de-blobbing story in full; in outline it does:

\begin{enumerate}
\item \textbf{Wire the octal pins.} \texttt{esp\_rom\_opiflash\_pin\_config} and
  the MSPI drive-strength setup connect the extra data lines and DQS that octal
  mode needs.
\item \textbf{Enable and size the PSRAM} via the from-source octal driver
  (\texttt{psram\_impl\_enable\_src} / \texttt{psram\_impl\_get\_physical\_size\_src}),
  reporting the detected size.
\item \textbf{Tune the data-in timing.} At 80\,MHz OPI-DDR the SPI0 data-in delay
  must be calibrated; \texttt{psram\_tune\_din} sweeps the candidate modes, measures
  which pass, and centres on the passing window (replacing an earlier hardcoded
  poke -- see Chapter~\ref{ch:psram}).
\item \textbf{Map it into the address space.} \texttt{Cache\_Dbus\_MMU\_Set} with
  the external-RAM flag maps the PSRAM at virtual base \texttt{16\#3D00\_0000\#}; the
  number of 64\,KB pages comes from the project's \texttt{board.ads}
  (\texttt{PSRAM\_Size / 64\,KB}).
\end{enumerate}

Because the bootloader maps PSRAM \emph{before} the application's \texttt{start.S}
sets up its stack, the application can even place the environment task's primary
stack in PSRAM (the \texttt{ENV\_STACK\_PSRAM} build option), trading a little
latency for a great deal of headroom on deep elaboration or recursion -- which
is exactly what the full-profile ACATS sweep needs.

\cnote{In ESP-IDF, \texttt{CONFIG\_SPIRAM=y} hides all of this. Here it is a few
dozen lines you can read: the pins, the enable call, the one mode-register fix,
and the MMU map. The size comes from \texttt{board.ads}, so a different module is
a one-line change, not a menuconfig safari.}

\section{\texttt{start.S}: from reset state to Ada}\index{start.S}

Once the loader jumps in, the application's \texttt{start.S}
(\texttt{examples/common/bare/start.S}) turns the raw post-loader state into
something Ada can run on:

\begin{itemize}
\item \textbf{Vectors and processor state.} It sets \texttt{VECBASE} to our own
  vector table (so any fault from here uses our handlers) and sets \texttt{PS} to
  the windowed-ABI state the runtime expects.
\item \textbf{Sane register windows.} It rewrites \texttt{WINDOWSTART} to a single
  frame at the current \texttt{WINDOWBASE}, discarding the loader's stale call
  frames so the first window overflow/underflow does not fault. (Getting this
  ordering right was the cure for an early environment-task double fault.)
\item \textbf{FPU on.} It sets \texttt{CPENABLE} so the first floating-point
  instruction does not trap.
\item \textbf{Stack and \texttt{.bss}.} It loads the stack pointer from the
  linker symbol and zeroes \texttt{.bss}.
\item \textbf{The instruction cache.} The loader enabled caching but left the
  I-cache unconfigured; \texttt{start.S} configures a 16\,KB, 8-way, 32-byte-line
  I-cache and sets the IROM/DROM MMU split, so instruction fetches from the IROM
  window actually return code rather than zeros.
\item \textbf{The clock.} This is the big one. Out of reset the CPU runs at
  \emph{20\,MHz} (the crystal divided by two); the runtime's timer math assumes
  240\,MHz, so short delays would run twelve times too long. \texttt{start.S}
  selects the 480\,MHz PLL divided to 240\,MHz (registers
  \texttt{16\#600C\_0010\#} and \texttt{16\#600C\_0060\#}) and tells the ROM timing
  helpers about the new frequency, all from IRAM before the jump into flash code.
\end{itemize}

It then jumps (not calls) into the C/Ada glue, which sets up the environment task
via \texttt{\_\_gnat\_enter\_env} -- a small shim that makes the env-task body the
\emph{outermost} window frame, so the first context switch has nothing stale
above it to corrupt. On the second core, a parallel cold-start entry
(\texttt{core1\_start}) repeats the window/VECBASE/stack normalisation before the
SMP runtime releases it.

\section{The memory map}\index{memory map}

All of this is pinned down by the linker scripts under
\texttt{examples/common/bare/vendor/} (the same regions appear in the memory-map
Table~\ref{tab:memmap}). The ESP32-S3's internal SRAM is visible at two
addresses -- an instruction-bus (IRAM) view and a data-bus (DRAM) view of the
\emph{same} bytes -- which is why code and data segments land in different
windows of one physical memory:

\begin{center}
\begin{tabular}{l l l}
\toprule
\textbf{Region} & \textbf{Base} & \textbf{What lives there} \\ \midrule
IRAM code     & \texttt{16\#4037\_8000\#} & vectors + hot code copied to SRAM \\
IROM (flash)  & \texttt{16\#4200\_0000\#} & bulk code, execute-in-place via cache \\
DRAM          & \texttt{16\#3FC9\_0000\#} & data, \texttt{.bss}, heap \\
DROM (flash)  & \texttt{16\#3C00\_0000\#} & read-only data via cache \\
PSRAM         & \texttt{16\#3D00\_0000\#} & external RAM (mapped by the loader) \\
RTC slow mem  & \texttt{16\#5000\_0000\#} & retained across deep sleep \\
\bottomrule
\end{tabular}
\end{center}

Two practical consequences fall out of this map and recur elsewhere in the book.
First, executable code must live in (or be reached through) an instruction-bus
window: that is why the full profile relocates GCC's nested-subprogram
trampolines to the IRAM alias rather than leaving them on the data-view stack
(\S\ref{ch:runtime-build}). Second, PSRAM is reachable only after the loader maps
it, and it cannot be a DMA source -- which is why DMA-capable scratch buffers are
carved from internal SRAM even when the task stack itself lives in PSRAM (the SD
and filesystem drivers rely on this).
